\begin{scanpage}{0090}{73}
\chaptertitle{III}{DIFFRACTION BY A FINITE LENGTH CYLINDER}
The distinctive feature of this problem is that, in addition to the nonuniform part of the current on the cylinder's surface which is caused by the discontinuity, there also exists a nonuniform part of the current arising as a consequence of the smooth curve of the surface. This part of the current has the character of waves travelling over the cylindrical surface along geodesic lines [36] --- that is, along spirals on the cylinder. These waves, which as they move strike the edge of the cylinder, undergo diffraction and excite secondary surface currents. In turn, the nonuniform part of the current resulting from the discontinuity undergoes diffraction while being propagated over the cylindrical surface. It is clear that specific consideration of all these effects is a very complicated problem.

However, if all the linear dimensions of the cylinder are sufficiently large in comparison with the wavelength, these effects may be neglected when calculating the fringing field in many cases which are of practical interest. In particular, they may be neglected when calculating the field scattered in the direction toward the
\end{scanpage}
